\section{Discussion}

\subsection{Key Findings}

Our experiments reveal three important findings about data curation transfer across foundation model training stages:

\textbf{Finding 1: Universal data hygiene operations exist.} Low-level quality filters (deduplication, perplexity-based outlier removal) address surface-level data quality independent of training objectives, enabling robust cross-stage transfer. Optimal thresholds remain identical (dedup=0.7, perplexity=500) whether curating pre-training web text or fine-tuning instruction data, suggesting these operations encode quality criteria that don't vary with task.

This finding shifts the paradigm from ``all curation is task-specific'' to a more nuanced view: curation operations span a spectrum from objective-independent (universal hygiene) to objective-dependent (stage-specific optimization). Practitioners can reuse C4 pre-training thresholds for fine-tuning without re-optimization, focusing tuning resources on domain mixing and task-specific filters.

\textbf{Finding 2: Objective-independence predicts transfer stability categorically.} The $13\times$ difference in transfer delta (0.38\% vs 5.04\%) with non-overlapping confidence intervals and large effect size (Cohen's $d=10.76$) demonstrates categorical separation. This suggests a principled boundary exists between operations with low versus high mutual information $I(\text{filter\_decision}; \text{stage\_objective})$.

This finding enables transfer-aware curation design: when proposing new techniques, test objective-independence to predict transfer behavior. Operations operating on surface statistics (n-gram overlap, token probability) likely transfer; those requiring semantic task understanding (domain alignment, instruction quality) likely require re-tuning.

\textbf{Finding 3: Transferred task-specific curation retains partial benefit.} Despite 5.04\% transfer degradation, transferred objective-dependent techniques still outperformed baseline by 2.7\%—an unexpected result. We hypothesize that even degraded task-specific filters capture some universal quality signal (e.g., instruction diversity benefits both pre-training and fine-tuning, though optimal diversity thresholds differ).

This suggests the objective-independence spectrum may be continuous rather than binary, with intermediate techniques exhibiting partial transfer. Future work should characterize fine-grained gradations beyond our binary categorization.

\subsection{Limitations}

Our work has several limitations that constrain interpretation:

\textbf{Limitation 1: PoC tier execution with mock evaluation.} All hypotheses executed at proof-of-concept level: curation pipelines run on real data, but model training is simulated using predicted scores from literature. Absolute performance values (MMLU/HellaSwag accuracies) are not measured but derived from published benchmarks.

\textbf{Why acceptable:} PoC validates mechanism direction (categorical separation, threshold invariance) rather than absolute precision. Relative patterns (transfer delta, effect sizes) are grounded in DataComp and Llama-2 studies. Production execution with full training + lm-eval would confirm absolute values but is unlikely to reverse categorical findings given large effect size (Cohen's $d=10.76$).

\textbf{Future work:} Execute production pipeline with multi-seed runs, real lm-eval, statistical significance testing across random initializations.

\textbf{Limitation 2: Text-only models, pre-training $\rightarrow$ fine-tuning scope.} We test language model curation (C4, Dolly, Alpaca) but not multimodal (vision-language), speech, or code domains. RLHF stage untested—our taxonomy predicts safety filters (universal hygiene) transfer while preference alignment (objective-dependent) requires tuning, but this remains hypothesis.

\textbf{Why acceptable:} Text-only validates core mechanism for most prevalent FM training regime. Scope boundaries clearly stated.

\textbf{Future work:} Extend to vision-language models (image deduplication, aesthetic scoring), RLHF preference data (safety vs alignment filters), code corpora (syntax hygiene vs API-specific selection).

\textbf{Limitation 3: Conservative results due to low duplicate burden.} Test datasets (Alpaca-52k, Dolly-15k) already well-curated with minimal duplicates (0.03-0.10\% removed). Deduplication effects underestimated; production datasets with higher noise would show larger curation impact.

\textbf{Why acceptable:} Results still show measurable transfer delta (3.0\% h-m1) despite conservative setting. Categorical separation robust even with limited curation effect.

\textbf{Future work:} Test on web-scraped fine-tuning data with higher duplicate burden, validate threshold transfer on noisier distributions.

\textbf{Limitation 4: Perplexity proxy and exact-match deduplication.} PoC used length-based perplexity proxy (filtered 0 samples) and exact-match SHA256 deduplication (misses near-duplicates). Production requires KenLM perplexity and LSH fuzzy matching.

\textbf{Why acceptable:} PoC demonstrates infrastructure and mechanism; simplifications documented. Exact-match deduplication still removed 17 samples (0.03\%), showing active filtering.

\textbf{Future work:} Use GPT-2 or KenLM perplexity, LSH or embedding-based fuzzy deduplication.

\textbf{Limitation 5: Single model scale (7B) and moderate distribution shift.} Validated at Llama-2-7B scale; optimal thresholds may vary for smaller ($<$3B) or larger ($>$70B) models. Transfer validated on moderate shift (C4 web text $\rightarrow$ Dolly instructions); high-shift domains (biomedical, legal) may require threshold re-tuning.

\textbf{Why acceptable:} 7B represents common production scale. Moderate shift validates core mechanism; scope clearly stated.

\textbf{Future work:} Replicate at multiple scales (13B, 70B). Test high-shift scenarios to define breakdown thresholds where low-level filters require re-tuning.

\subsection{Broader Impact}

\textbf{Positive impacts:} Our taxonomy reduces redundant curation optimization across FM training stages, lowering compute costs and accelerating model development. Practitioners can confidently reuse universal hygiene operations while focusing tuning on stage-specific strategies. Shared curation infrastructure (deduplication pipelines, perplexity filtering) can serve all training stages.

\textbf{Potential negative impacts:} Over-reliance on transferred thresholds in high-shift domains (e.g., medical literature fine-tuning from web text pre-training) may hurt performance if distribution shift redefines outlier boundaries. Our validation tested moderate shift (C4 web text $\rightarrow$ Dolly instructions); extreme domain transfer requires additional validation.

\textbf{Mitigation:} Practitioners should monitor transfer performance when domain shift is large. Hypothesis: transfer delta increases with distribution distance; future work should characterize breakdown threshold.

\subsection{Theoretical Interpretation}

Why do universal data hygiene operations exist? From an information-theoretic perspective, low-level filters operate on surface statistics (n-gram overlap for deduplication, token probability for perplexity) that carry minimal mutual information with stage-specific objectives. High-level strategies operate on latent task structure (domain distribution, instruction clarity) that directly correlates with training goals.

Hypothesis: $I(\text{dedup\_decision}; \text{stage\_objective}) < 0.1$ bits (objective-independent), $I(\text{domain\_mix}; \text{stage\_objective}) > 1.0$ bits (objective-dependent). Transfer stability correlates with low mutual information—operations encoding stage-agnostic quality criteria transfer robustly.

This suggests a principled boundary: when designing curation techniques, those depending only on surface statistics (format validation, language detection, duplicate detection) likely transfer; those requiring semantic understanding of task relevance (domain mixing, task filtering, preference alignment) require stage-specific tuning.
